% !TeX program = lualatex
% !TeX encoding = utf8
% !TeX spellcheck = uk_UA
% !TeX root =./ClassicalElectrodynamics.tex

%=========================================================
\chapter*{Вибір системи одиниць}\label{\currfilebase}
\addcontentsline{toc}{chapter}{Вибір системи одиниць}
%=========================================================

У класичній механіці природно будувати систему одиниць як абсолютну, тобто таку, в якій усі фізичні величини виражаються через базові --- довжину, масу та час. Однією з таких систем є метрична система СГС, що базується на трьох основних одиницях: сантиметрі (довжина), грамі (маса) та секунді (час), --- і власне назва <<СГС>> утворена саме від перших літер цих одиниць.

Цей підхід можна послідовно поширити й на електродинаміку, якщо розглядати електричні та магнітні величини як похідні від механічних. Таке поширення реалізоване в гаусовій системі одиниць --- різновиді системи СГС. У цій системі електричні величини вимірюються в електростатичних одиницях (СГСЕ), а магнітні --- в електромагнітних (СГСМ). Одиницю заряду тут визначають безпосередньо через сантиметр, грам і секунду за допомогою закону Кулона \(F = \dfrac{q_1 q_2}{r^2}\).

На противагу природному механічному базису, у SI штучно введено ще одну фундаментальну одиницю --- ампер, яка не має прямого фізичного змісту, адже електричний струм є лише похідною величиною від руху зарядів. Одиниця заряду в SI (кулон) є похідною від ампера, що вивертає логічний зв'язок: замість того щоб визначати струм через заряд, заряд визначають через струм. Хоча система SI є зручною для практичних та інженерних застосувань, оскільки результати обчислень одразу отримуються у звичних одиницях --- вольтах, амперах, омах тощо, натомість, гаусова система дозволяє зосередитися на фізичному змісті явищ, не обтяжуючи формули зайвими коефіцієнтами та допоміжними величинами, в цій системі рівняння не містять діелектричної \( \varepsilon _{0}\) чи магнітної \(\mu _{0}\) сталих, які не мають фізичного змісту, електричні $\Efield$ та магнітні $\Bfield$ поля мають однакову розмірність, зв'язок одиниць легко відслідковується з рівнянь. Все це робить формули теорії <<елегантнішими>> та симетричнішими, у результаті чого теорія набуває більш цілісного та компактного вигляду. Гаусова система залишається стандартом у теоретичній електродинаміці та багатьох галузях сучасної фізики.

З огляду на це, у подальшому викладі всі основні результати формулюються в гаусовій системі. За потреби перехід до одиниць SI може бути виконаний наприкінці обчислень.
